% Integrantes:
% - Arthur Gama Ruiz (RA: 23.01445-8)
% - Enzo Oliveira D’Onofrio (RA: 23.01561-6)
% - Felipe Fazio da Costa (RA: 23.00055-4)
% - João Vitor Morimoto Sesma (RA: 23.01516-0)
% - Leonardo Souza Olivieri (RA: 23.01512-8)
% - Pedro Wilian Palumbo Bevilacqua (RA: 23.01307-9)
%
% Data: 04/06/2026
% Matérias: 
% - ECM516_Arquitetura_de_Computadores
% - ECM252_Linguagens_de_Programação_2

\documentclass[12pt,a4paper]{article}

% --- Pacotes Básicos ---
\usepackage[utf8]{inputenc}
\usepackage[T1]{fontenc}
\usepackage[portuguese]{babel}
\usepackage{amsmath, amsfonts, amssymb}
\usepackage{graphicx}
\usepackage{geometry}
\usepackage{setspace}
\usepackage{indentfirst}
\usepackage{titlesec}
\usepackage{tikz}
\usepackage{listings}
\usepackage{xcolor}
\usepackage{hyperref}
\usepackage{caption}

% --- Configurações de Margens (ABNT) ---
\geometry{
    top=3cm,
    left=3cm,
    right=2cm,
    bottom=2cm
}

% --- Configurações de Espaçamento e Fonte ---
\onehalfspacing 
\setlength{\parindent}{1.25cm} 

% --- Estilo de Código ---
\definecolor{codegreen}{rgb}{0,0.6,0}
\definecolor{codegray}{rgb}{0.5,0.5,0.5}
\definecolor{codepurple}{rgb}{0.58,0,0.82}
\definecolor{backcolour}{rgb}{0.95,0.95,0.92}

\lstdefinestyle{mystyle}{
    backgroundcolor=\color{backcolour},   
    commentstyle=\color{codegreen},
    keywordstyle=\color{magenta},
    numberstyle=\tiny\color{codegray},
    stringstyle=\color{codepurple},
    basicstyle=\ttfamily\footnotesize,
    breakatwhitespace=false,         
    breaklines=true,                 
    captionpos=b,                    
    keepspaces=true,                 
    numbers=left,                    
    numbersep=5pt,                  
    showspaces=false,                
    showstringspaces=false,
    showtabs=false,                  
    tabsize=2
}
\lstset{style=mystyle}

% --- Definição da Linguagem JavaScript ---
\lstdefinelanguage{JavaScript}{
  morekeywords={async, await, break, case, catch, class, const, continue, debugger, default, delete, do, else, enum, export, extends, false, finally, for, function, if, import, in, instanceof, let, new, null, return, super, switch, this, throw, true, try, typeof, var, void, while, with, yield},
  sensitive=true,
  morecomment=[l]{//},
  morecomment=[s]{/*}{*/},
  morestring=[b]',
  morestring=[b]"
}

% --- TikZ Libraries ---
\usetikzlibrary{shapes,arrows,positioning,calc}

% --- Início do Documento ---
\begin{document}

% --- CAPA ---
\begin{titlepage}
    \begin{center}
        {\large \textbf{CEUN-IMT - INSTITUTO MAUÁ DE TECNOLOGIA}} \\
        \vspace{0.5cm}
        {\large CURSO DE ENGENHARIA DE COMPUTAÇÃO} \\
        \vspace{0.5cm}
        {\large ECM516 - ARQUITETURA DE COMPUTADORES} \\
        {\large ECM252 - LINGUAGENS DE PROGRAMAÇÃO 2} \\
        \vfill
        {\Large \textbf{SkySwift: Desenvolvimento de uma Infraestrutura de Entrega Autônoma via Drones baseada em Microsserviços}} \\
        \vfill
        {\large Arthur Gama Ruiz - 23.01445-8} \\
        {\large Enzo Oliveira D’Onofrio - 23.01561-6} \\
        {\large Felipe Fazio da Costa - 23.00055-4} \\
        {\large João Vitor Morimoto Sesma - 23.01516-0} \\
        {\large Leonardo Souza Olivieri - 23.01512-8} \\
        {\large Pedro Wilian Palumbo Bevilacqua - 23.01307-9} \\
        \vfill
        {\large SÃO CAETANO DO SUL} \\
        {\large 2026}
    \end{center}
\end{titlepage}

% --- FOLHA DE ROSTO ---
\newpage
\begin{center}
    {\large \textbf{Arthur Gama Ruiz, Enzo Oliveira D’Onofrio, Felipe Fazio da Costa, \\ João Vitor Morimoto Sesma, Leonardo Souza Olivieri, Pedro Bevilacqua}} \\
    \vfill
    {\Large \textbf{SkySwift: Desenvolvimento de uma Infraestrutura de Entrega Autônoma via Drones baseada em Microsserviços}} \\
    \vfill
    \hspace{.45\textwidth}
    \begin{minipage}{.5\textwidth}
        Relatório técnico apresentado como requisito parcial para aprovação nas disciplinas de Arquitetura de Computadores (ECM516) e Linguagens de Programação 2 (ECM252) no CEUN-IMT - Instituto Mauá de Tecnologia.
    \end{minipage}
    \vfill
    {\large SÃO CAETANO DO SUL} \\
    {\large 2026}
\end{center}

% --- RESUMO ---
\newpage
\begin{center}
    {\large \textbf{RESUMO}}
\end{center}
\vspace{1cm}
Este trabalho apresenta o projeto SkySwift, uma solução integrada para logística autônoma utilizando drones. A arquitetura proposta rompe com o modelo monolítico tradicional, adotando um ecossistema de microsserviços desacoplados que se comunicam através de um barramento de eventos centralizado. A implementação técnica abrange desde o backend em Node.js com persistência poliglota (MySQL e SQLite), até uma interface de usuário rica desenvolvida em React.js. O foco principal reside na confiabilidade da entrega de mensagens, segurança através de criptografia assimétrica e JWT, e precisão geográfica via algoritmos de roteirização. Os resultados demonstram um sistema resiliente, capaz de processar pedidos, calcular rotas e notificar usuários de forma assíncrona e escalável.

\textbf{Palavras-chave:} Microsserviços, Barramento de Eventos, Node.js, React, Logística de Drones.

% --- LISTAS ---
\newpage
\tableofcontents
\newpage
\listoffigures

% --- INTRODUÇÃO ---
\newpage
\section{INTRODUÇÃO}
O mercado de entregas \textit{last-mile} enfrenta desafios crescentes relacionados à infraestrutura urbana. A proposta do SkySwift é automatizar esse processo utilizando drones, integrando tecnologias de ponta em desenvolvimento de software e arquitetura de sistemas distribuídos. Este relatório detalha a engenharia por trás da plataforma, focando na separação de responsabilidades e na robustez da comunicação entre os módulos.

\section{ARQUITETURA DE MICROSSERVIÇOS}

\subsection{Filosofia de Design}
Diferente de uma aplicação monolítica, onde uma falha em um módulo pode comprometer todo o sistema, o SkySwift utiliza o isolamento de processos. Cada serviço possui seu próprio contexto, dependências e, em alguns casos, banco de dados.

\subsection{Barramento de Eventos (Event Bus)}
O coração da arquitetura é o \textbf{Barramento de Eventos}. Ele atua como um mediador (\textit{Middleware}) seguindo o padrão \textit{Observer}.

\textbf{Funcionamento:}
1. Um serviço de origem realiza um \textit{POST} para o barramento enviando um objeto JSON que contém o \texttt{tipo} do evento e os \texttt{dados}.
2. O barramento registra o evento em um histórico.
3. O barramento percorre uma lista interna de serviços "inscritos" e repassa o evento para cada um deles via HTTP.

\begin{figure}[h]
    \centering
    \begin{tikzpicture}[node distance=2cm, auto]
        \tikzstyle{svc} = [rectangle, draw, fill=orange!20, text width=5em, text centered, rounded corners, minimum height=2em]
        \tikzstyle{bus} = [rectangle, draw, fill=blue!30, text width=8em, text centered, rounded corners, minimum height=4em]
        
        \node [svc] (source) {Serviço Origem};
        \node [bus, right=2cm of source] (bus) {BARRAMENTO\\(Hub Central)};
        \node [svc, right=2cm of bus, yshift=1.5cm] (sub1) {Inscrito A};
        \node [svc, right=2cm of bus, yshift=0cm] (sub2) {Inscrito B};
        \node [svc, right=2cm of bus, yshift=-1.5cm] (sub3) {Inscrito C};

        \draw [->, thick] (source) -- node[above, font=\tiny]{Publica Evento} (bus);
        \draw [->, thick] (bus) -- (sub1);
        \draw [->, thick] (bus) -- (sub2);
        \draw [->, thick] (bus) -- (sub3);
    \end{tikzpicture}
    \caption{Fluxo de Distribuição de Eventos (Pub/Sub)}
\end{figure}

\newpage
\section{DETALHAMENTO TÉCNICO E CÓDIGO}

\subsection{Barramento de Eventos (Hub Central)}
O Barramento de Eventos é o sistema nervoso central do ecossistema. Ele gerencia a inscrição de microsserviços e a distribuição de mensagens seguindo o modelo \textit{Publish/Subscribe}. Nesta arquitetura, o serviço emissor (Publisher) não precisa conhecer o endereço ou a existência dos receptores (Subscribers), o que reduz a complexidade de conexão de $O(n^2)$ para $O(n)$.

A implementação do broadcast no arquivo \texttt{server.js} do barramento é crítica para a consistência eventual do sistema. Abaixo, o trecho que garante a entrega para todos os serviços ativos:

\begin{lstlisting}[language=JavaScript, caption=Implementação de Distribuição de Eventos]
async function distribuirEvento(evento) {
  const resultados = [];
  for (const servico of inscricoes) {
    if (servico.nome === evento.origem) continue; // Evita loop infinito

    try {
      await axios.post(`${servico.url}/eventos/receber`, evento, {
        timeout: 5000,
      });
      resultados.push({ servico: servico.nome, status: 'entregue' });
    } catch (erro) {
      resultados.push({ servico: servico.nome, status: 'falha' });
    }
  }
  return resultados;
}
\end{lstlisting}

\subsection{Cálculo de Rota e Rastreamento}
O microsserviço \texttt{entrega\_via\_drone} implementa uma estratégia de \textit{fallback} geográfica. Ele tenta consumir o provedor OSRM (\textit{Open Source Routing Machine}), mas possui uma inteligência local baseada na Fórmula de Haversine. Esta abordagem garante que, mesmo sem conectividade com a malha rodoviária externa, o drone consiga estimar o tempo de voo em linha reta.

\begin{equation}
d = 2R \cdot \operatorname{atan2}\left(\sqrt{a}, \sqrt{1-a}\right)
\end{equation}
Onde $a = \sin^2(\Delta\phi/2) + \cos \phi_1 \cdot \cos \phi_2 \cdot \sin^2(\Delta\lambda/2)$.

\begin{lstlisting}[language=JavaScript, caption=Implementação da Fórmula de Haversine]
function haversineDistanceMeters(lat1, lon1, lat2, lon2) {
  const R = 6371000; // Raio da Terra em metros
  const dLat = (lat2 - lat1) * Math.PI / 180;
  const dLon = (lon2 - lon1) * Math.PI / 180;
  const a = Math.sin(dLat/2) * Math.sin(dLat/2) +
            Math.cos(lat1 * Math.PI / 180) * Math.cos(lat2 * Math.PI / 180) *
            Math.sin(dLon/2) * Math.sin(dLon/2);
  const c = 2 * Math.atan2(Math.sqrt(a), Math.sqrt(1-a));
  return R * c;
}
\end{lstlisting}

\subsection{Frontend React e Orquestração de Estado}
A interface foi desenvolvida utilizando \textbf{React 19} com \textbf{Vite}. O grande diferencial técnico reside no roteamento customizado baseado em estado, que evita o uso de bibliotecas externas pesadas e garante transições rápidas entre o painel de pedidos e o mapa de rastreamento.

\begin{figure}[h]
    \centering
    \begin{tikzpicture}[node distance=1.5cm, auto]
        \node [rectangle, draw, fill=blue!10] (state) {Estado Global (Auth/Tema)};
        \node [rectangle, draw, below left=of state] (page1) {Página Pedidos};
        \node [rectangle, draw, below=of state] (page2) {Rastreamento};
        \node [rectangle, draw, below right=of state] (page3) {Perfil};
        \draw [->, thick] (state) -- (page1);
        \draw [->, thick] (state) -- (page2);
        \draw [->, thick] (state) -- (page3);
    \end{tikzpicture}
    \caption{Fluxo de Estado e Propagação de Dados no Frontend}
\end{figure}

\subsection{Segurança e Integração Reativa}
O serviço de usuários protege os dados sensíveis utilizando \texttt{bcrypt}. A autenticação gera um \textit{JSON Web Token} (JWT) que permite o processamento distribuído (Arquitetura de Computadores), onde cada microsserviço pode validar a identidade do usuário sem sobrecarregar o banco de dados central, utilizando apenas a chave pública de assinatura.

Serviços periféricos como \texttt{notificacoes} e \texttt{contato\_email} funcionam como consumidores reativos puros. Eles processam eventos do tipo \texttt{PEDIDO\_CRIADO} para disparar e-mails via SMTP e alertas em tempo real.

\begin{figure}[h]
    \centering
    \begin{tikzpicture}[node distance=2.5cm, auto]
        \node [circle, draw, fill=yellow!20] (event) {Evento};
        \node [rectangle, draw, right=of event, fill=green!10] (logic) {Processamento};
        \node [rectangle, draw, right=of logic, fill=orange!10] (action) {Ação (SMTP/Push)};
        \draw [->, thick] (event) -- (logic);
        \draw [->, thick] (logic) -- (action);
    \end{tikzpicture}
    \caption{Cadeia de Reação Reativa de Serviços Periféricos}
\end{figure}

\section{INFRAESTRUTURA DE DADOS}
O sistema utiliza um esquema relacional para garantir a integridade referencial. Abaixo, o diagrama simplificado das tabelas principais:

\begin{figure}[h]
    \centering
    \begin{tikzpicture}[node distance=3cm, every node/.style={rectangle, draw, fill=gray!10, text width=8em, align=left}]
        \node (user) {\textbf{USUARIOS}\\ id (PK)\\ nome\\ email\\ senha (hash)};
        \node (order) [right=of user] {\textbf{PEDIDOS}\\ id (PK)\\ user\_id (FK)\\ status\\ destino\_lat\\ destino\_lng};
        \node (notify) [below=1.5cm of order] {\textbf{NOTIFICACOES}\\ id (PK)\\ pedido\_id (FK)\\ mensagem\\ lida (bool)};

        \draw [->, thick] (user) -- (order);
        \draw [->, thick] (order) -- (notify);
    \end{tikzpicture}
    \caption{Modelo de Dados Relacional}
\end{figure}

\section{CONCLUSÃO}
O SkySwift demonstrou ser uma plataforma robusta para a gestão de entregas autônomas. A escolha pela arquitetura de microsserviços via barramento de eventos provou-se eficaz para o desacoplamento, permitindo que novos serviços (como um gateway de pagamento) sejam adicionados sem alterar o código existente. O projeto atende integralmente aos requisitos de ECM516 e ECM252, consolidando conhecimentos de sistemas distribuídos, segurança e desenvolvimento frontend moderno no CEUN-IMT.

\end{document}
